\chapter{1938:Richard Kuhn}

\label{chap:chemistry1938}

The Nobel Prize in Chemistry for the year 1938 was awarded to Richard Kuhn.

\textbf{Motivation:}

for his work on carotenoids and vitamins

\section{Richard Kuhn}

\subsection*{Early Life and Education}

Richard Kuhn was born on 1900-12-03 in Vienna, Austria-Hungary (now Austria).
He was an Austrian-German biochemist awarded the Nobel Prize in Chemistry in 1938 for his work on carotenoids and vitamins.

\subsection*{Career and Major Research}

Kuhn studied chemistry in Vienna and Munich and received his doctorate at the University of Munich in 1922 under Richard Willstätter.
He was professor at the Eidgenössische Technische Hochschule in Zürich from 1926 to 1929.
In 1929 he became director of the Institute of Chemistry of the Kaiser Wilhelm Institute for Medical Research in Heidelberg, where he worked for the rest of his life.
He investigated the carotenoid pigments, determining the structures of many of them, and his laboratory prepared vitamin A from beta-carotene by careful purification.
He also isolated riboflavin, the vitamin known as lactoflavin (vitamin B2), established its structure, and achieved its synthesis.
In 1938 he was awarded the Nobel Prize in Chemistry, but the German government forbade him to accept it; he received the prize only in 1949.

\subsection*{Nobel Prize-winning Work}

The work for which Richard Kuhn was awarded the Nobel Prize in Chemistry in 1938 was described by the Nobel Committee as: "for his work on carotenoids and vitamins".

Kuhn's prize recognized his work on the carotenoids and vitamins.
He established the constitutions of many carotenoid pigments and of riboflavin (vitamin B2), and he achieved the synthesis of riboflavin, confirming its structure.

\subsection*{Significance and Impact}

Kuhn's investigations laid the chemical foundations for the understanding of the carotenoids and of the B vitamins.
His synthesis of riboflavin made the vitamin available in quantity and opened the way to studying its biological function as a component of the flavin coenzymes.

\subsection*{Later Career and Honors}

Kuhn continued to direct his institute in Heidelberg after the war, and his later research dealt with enzymes and with natural substances such as the antibiotics.
He held honorary doctorates from several universities.
He died on 1967-08-01 in Heidelberg, Germany.

\subsection*{Legacy}

The legacy of Richard Kuhn endures through the lasting impact of his scientific contributions.
The chemistry of the carotenoids and vitamins that he established continues to be built upon by researchers across the chemical and biological sciences.

\subsection*{Modern Developments}

Kuhn's synthesis and structural work on vitamins, including vitamin B2 and vitamin B6, helped launch the modern vitamin industry and the understanding of coenzymes. His investigations of carotenoids and natural products continue to inform pigment chemistry and nutrition science. Kuhn was also a pioneer in quantum chemistry before focusing on vitamins.

\subsection*{Selected Publications}

\begin{itemize}

\item Kuhn, R., \& Weygand, F. (1934). "Die Synthese des Lactoflavins." Berichte der deutschen chemischen Gesellschaft.

\end{itemize}

Kuhn's extensive experimental work on the carotenoids and vitamins was published in numerous papers, most of them in Berichte der deutschen chemischen Gesellschaft and other German journals of the period.

\clearpage

\section*{Chapter Summary}

The 1938 prize recognized Richard Kuhn for his work on carotenoids and vitamins. He established the structures of riboflavin (vitamin B2) and several carotenoids, including the chemistry of the vitamin B group.
